\documentclass[12pt, a4paper]{article}

\usepackage[utf8]{inputenc}
\usepackage[T1]{fontenc}

\DeclareUnicodeCharacter{00B0}{\ensuremath{^\circ}}
\usepackage{amsmath, amssymb}
\usepackage{graphicx}
\usepackage{xcolor}
\usepackage{booktabs}
\usepackage{natbib}
\usepackage{hyperref}
\usepackage[margin=2.5cm]{geometry}
\usepackage{lineno}
\usepackage{setspace}
\usepackage{siunitx}
\usepackage{cleveref}
\usepackage{placeins}
\usepackage{rotating}

\doublespacing
\linenumbers


\newcommand{\todo}[1]{\textcolor{red}{\textbf{[TODO: #1]}}}
\newcommand{\placeholder}[1]{\textcolor{gray}{\textit{#1}}}
\newcommand{\Sv}{\,\text{Sv}}    % 1 Sv = 10^6 m^3/s
\newcommand{\psu}{\,\text{psu}}
\newcommand{\CO}{CO\ensuremath{_2}}
\newcommand{\dint}{\mathrm{d}}


\begin{document}

\title{Coherence of Resilience Metrics for the Atlantic Meridional
       Overturning Circulation Across a Model Hierarchy}

\author{Andreas Morr$^{1,2}$%
\thanks{Corresponding author: \texttt{andreas.morr@tum.de}}
\and
Reyk Börner$^{3}$
\and
Sjoerd Terpstra$^{3}$
\and
Christian Kuehn$^{1}$}


\date{\today}
\maketitle

\begin{abstract}
The Atlantic Meridional Overturning Circulation (AMOC) is a critical component of the global climate system whose potential collapse constitutes one of the most consequential climate tipping points. Quantifying AMOC resilience, that is, its capacity to maintain the active circulation state under external perturbations, is therefore of high scientific and societal importance. Resilience is a multifaceted concept, encompassing notions such as local and global recovery times, basin volume, and the minimal perturbation required to trigger a collapse. Whether these distinct metrics yield consistent, and therefore trustworthy, signals across models of different complexity has not been systematically examined. Here, we address this question by analysing four models that span a hierarchy in physical complexity: (i) a 3-box salinity model, parameterised using comprehensive model output; (ii) a two-dimensional latitude-depth Boussinesq model of the overturning circulation; (iii) the Earth system model of intermediate complexity CLIMBER-X; and (iv) the general circulation model PlaSim-LSG. Across the hierarchy, we compute model-specific implementations of four complementary resilience concepts: local resilience, convergence time, basin volume, and minimal critical shock. Our results suggest that not all metrics are coherent across the model hierarchy, as exhibited by incompatible trends under increasing \CO{} forcing in each model. Several measures exhibit an unexpected increase of system resilience under climate change. Overall coherence of resilience trends is limited, pointing to deficiencies in model compatibility and challenges in the application of these concepts in real-world systems These findings have direct implications for the detectability of an imminent AMOC collapse from observational or model-based early-warning signals.

\end{abstract}

\textbf{Keywords:} Atlantic Meridional Overturning Circulation, tipping points, resilience, model hierarchy

\newpage

\section{Introduction}
\label{sec:intro}

The Atlantic Meridional Overturning Circulation is one of the dominant features of the global ocean, transporting large amounts of warm, saline surface waters northward across the Atlantic and returning cold, deep North Atlantic water southward at depth. This large-scale transport pattern underpins the relatively mild climate of Western Europe and plays a central role in the global carbon and heat cycles. Palaeoclimate records document repeated rapid transitions between an active (AMOC-on) and a weakened or collapsed (AMOC-off) state during the Pleistocene, associated with Dansgaard–Oeschger events. Under continued anthropogenic forcing, rising sea surface temperatures, freshening of the North Atlantic from increased precipitation, and Greenland ice-sheet melt may push the AMOC towards a similar transition.

More concretely, in background climate configurations where at least two stable equilibria coexist (the on and off states), the system can be permanently tipped from one to the other either by crossing a critical threshold leading to the annihilation of one of the formerly stable states (bifurcation-induced tipping) or by a sufficiently large state perturbation (shock-induced tipping) \citep{Ashwin2012tipTypes, Alkhayuon2019AMOCRateTipping}. In the bifurcation-theory framework for AMOC tipping, the transition to the off state is often conceptualised as a saddle-node bifurcation \citep{VanWesten2025SaddleNodeAMOC}. A shock-induced transition, on the other hand, can be understood as an exit of the system state from the basin of attraction associated with the AMOC on state. In both frameworks, the AMOC is investigated as an autonomous dynamical system, potentially influenced by slow parameter changes or external perturbations, such as noise. This allows for the characterisation of the system in terms of several quantities, subsumed under the term "resilience" \citep{Krakovska2024ResilienceDynamicalSystems}.

One important notion of resilience has been applied to the AMOC in the context of a bifurcation-induced transition. As the system approaches such a transition through a gradual change in the background climate, the stable on-state loses linear stability, a phenomenon known as critical slowing down (CSD) \citep{Scheffer2009EWSNature}. CSD manifests in observable statistics: near the bifurcation, the dominant eigenvalue of the linearised system approaches zero, leading to increases in variance and autocorrelation of system fluctuations \citep{Held2004CSDDetection, Dakos2008CSDClimate}. These generic signatures have motivated a large body of research on early-warning indicators of tipping \citep{Lenton2011EarlyWarningClimate, Boers2021AMOCEWS, Ditlevsen2023AMOCPrediction}.

Yet CSD-based indicators capture only the system's reaction to small perturbations around the equilibrium, known as local resilience, quantifying how the on-state changes locally as the bifurcation is approached. A complementary perspective is provided by non-local resilience concepts, which characterise the size and geometry of the basin of attraction of the on-state and therefore also capture the system's capacity to recover from larger perturbations \citep{Holling1973ResilienceStabilityEcological, Gunderson2000EcologicalResilienceTheory, Menck2013BasinStability}. A comprehensive assessment of the resilience of the AMOC on state should therefore encompass many different notions. In this work, we consider four such choices. The following heuristic definitions explain the interpretation of the respective resilience measure. The concrete implementation and calculation in each model will vary and are explained in Section \ref{sec:methods}.

\begin{enumerate}
  \item \textbf{Local resilience}: capacity of a system to continuously absorb small perturbations to its equilibium on-state, such as noise and internal variability.
  \item \textbf{Convergence time}: duration trajectories take to return to the on-state equilibrium after a large perturbation.
  \item \textbf{Basin volume}: the volume fraction of a reference region of initial model conditions that converges to the on-state \citep{Menck2013BasinStability}.
  \item \textbf{Minimal critical shock}: the smallest instantaneous perturbation that moves the system out of the basin of attraction of the on-state \citep{Klinshov2015DistanceThresholdNumerical, Halekotte2020MinimalFatalShocks}.
\end{enumerate}

In natural systems, controlled perturbation experiments are rarely feasible, making direct estimation of non-local resilience quantities from observations difficult. This is especially true for the AMOC, and numerical models, therefore, play an essential role. However, climate models span an enormous range of complexity, from analytically tractable conceptual models to full Earth system models (ESMs), and it is not obvious whether resilience metrics computed in simple models translate to their more complex counterparts, which include many additional feedbacks and degrees of freedom \citep{Stouffer2006InvestigatingCausesAMOC, Schneider2017EarthSystemModeling}. If metrics disagree across the hierarchy, it could mean that their utility as real-world system characteristics is severely limited, or that parts of the hierarchy do not expose them well. Several studies have examined individual resilience concepts in the context of the AMOC at specific levels of model complexity \citep{Rahmstorf2005AMOCHystersis, Cimatoribus2014AMOCBoxModel, Boers2021AMOCEWS, Ditlevsen2023AMOCPrediction, Jacques-Dumas2024RareEventResilience}, but a systematic, multi-metric comparison across a model hierarchy is lacking.

Our analysis combines four modelling studies that span many orders of magnitude of model complexity: (i) a conceptual 3-box salinity model \citep{Wood20195BoxModel, Alkhayuon2019AMOCRateTipping}; (ii) a latitude-depth Boussinesq partial differential equation model \citep{Esclapez2026ConstructingEfficientScore}; (iii) the Earth system model of intermediate complexity CLIMBER-X \citep{Willeit2022EarthSystemModel}; and (iv) the coupled general circulation model PlaSim-LSG \citep{Borner2025AMOCGlobalStability}. Section~\ref{sec:methods} describes each model and the resilience metrics computed for it. Section~\ref{sec:results} reports the individual model's results and Section~\ref{sec:discussion} synthesises them across the hierarchy. In Section~\ref{sec:outlook}, we discuss implications for research on the AMOC and on early-warning detection of potential tipping.

\section{Methods}
\label{sec:methods}

We compute a common set of resilience quantities for each model while respecting the different experimental possibilities at each level of the hierarchy. For this, we analyse trajectories from perturbed initial conditions in each model. The reduced state space is defined by North Atlantic and Tropical Atlantic salinity values. In the Boussinesq and CLIMBER-X experiments, salt is conserved by applying a compensating perturbation in a Southern box; in the 3-box model, salt conservation is inherent to the model equations via a dependent Indo-Pacific salinity variable . Figure \ref{fig:perturbations-overview} illustrates the model-specific perturbation and diagnostic regions, and their implementation is explained in the respective subsections below. Table \ref{tab:overview} gives an overview of the concrete meaning and estimation routine for each resilience measure across the models. Initial conditions in the 3-box, Boussinesq, and CLIMBER-X experiments are probed within rectangles bounded by $\pm2\psu{}$ perturbations relative to their model-specific reference configurations. This is done by drawing a large amount of uniform samples from this region (3-box model) or by evenly sampling a grid of initial conditions (Boussinesq model and CLIMBER-X). PlaSim-LSG constitutes an exception: we analyse existing edge-track and equilibrium runs from \cite{Borner2025AMOCGlobalStability} using two salinity readouts in the North and Tropical Atlantic.

Within each model, we repeat the analysis across a global warming forcing coordinate to be specified for each model below. This allows us to compare directional trends of resilience, which are more meaningful across models than absolute values expressed in model-specific units. The resilience of the AMOC on-state is the primary object of interest throughout, though all analyses could also be carried out for the off-state whenever it is present.





\begin{figure}[t]
  \centering
  \includegraphics[width=\textwidth,height=0.82\textheight,keepaspectratio]{perturbations_overview.png}
  \caption{\textbf{Perturbation and diagnostic geometry across the model hierarchy.}
  This overview compares the geographic salinity regions used as perturbation coordinates (shaded regions) in the Boussinesq, and CLIMBER-X experiments and as diagnostic readouts (dashed boxes) in the Boussinesq, and CLIMBER-X, and PlaSim-LSG analysis. The box geometry associated with the 3-box model is approximately reproduced from the model conception in \cite{Wood20195BoxModel}.}
  \label{fig:perturbations-overview}
\end{figure}

\begin{table}[t]
  \centering
  \small
  \setlength{\tabcolsep}{4pt}
  \caption{\textbf{Overview of resilience-metric implementations across the model hierarchy.}}
  \label{tab:overview}
  \begin{tabular}{p{2.3cm} p{3.2cm} p{3.2cm} p{3.2cm} p{3.2cm}}
    \toprule
    & \textbf{3-Box model} & \textbf{Boussinesq} & \textbf{CLIMBER-X} & \textbf{PlaSim-LSG} \\
    \midrule
    Local resilience
      & $-\mathrm{Re}(\lambda_{\max})$: negative real part of the least stable eigenvalue of the Jacobian at the on-state fixed point, evaluated via forward-mode automatic differentiation \citep{Khan2015VectorForwardMode}.
      & $-\mathrm{Re}(\lambda_{\max})$: negative real part of the least stable eigenvalue fitted from a $2\times2$ discrete linear map via isotropic salinity perturbations.
      & $-\mathrm{Re}(\lambda_{\max})$: negative real part of the least stable eigenvalue fitted from a $2\times2$ discrete linear map via isotropic salinity perturbations.
      & Inverse of the long axis of the $1\sigma$ Gaussian ellipse in the leading (EOF1, EOF2) plane: $1/\ell_{\max}$, where $\ell_{\max} = 2000\sqrt{\lambda_{\max}(C_{12})}$ is the scaled semi-major axis length, $C_{12}$ is the $2\times2$ marginal covariance of the on-state equilibrium run, and $\lambda_{\max}$ denotes the largest eigenvalue \todo{check this}. \\[6pt]
    Convergence time
      & Mean return time of on-state trajectories among $10\,000$ uniform random initial conditions in a $\pm 2\psu{}$ box; convergence at $\varepsilon=0.05\psu{}$ proximity.
      & Mean return time of valid on-state trajectories on a $20\times20$ grid over $\pm 2\psu{}$; convergence at $\varepsilon=0.05\psu{}$ proximity.
      & Mean return time of completed trajectories on a $9\times9$ grid over $\pm 2\psu{}$; convergence defined by permanent entry into an $\varepsilon=0.05\psu{}$ ball.
      & Mean transit time from the last edge-state $4\sigma$ ellipse visit to the first on-state $4\sigma$ ellipse entry. \\[6pt]
    Basin volume
      & Fraction of $10\,000$ random initial conditions in a $\pm 2\psu{}$ box centred on the on state that converge to the on state.
      & Fraction of valid initial conditions on the $20\times20$ perturbation grid converging to the on state.
      & Not applicable \todo{rethink}.
      & Not estimated because the available edge tracks do not sample a reference region of state space. \\[6pt]
    Minimal critical shock
      & Minimum distance from the on state to a sampled initial conditions in a region of $\pm 2$psu around the on state not converging to the on state.
      & Minimum salinity-plane distance from the on state to a valid sampled initial condition not converging to the on state.
      & Not applicable: \todo{rethink}.
      & Gap between the on-state and edge-state $4\sigma$ Gaussian ellipses in the deep-box salinity plane. \\
    \bottomrule
  \end{tabular}
\end{table}


\subsection{3-Box salinity model}
\label{sec:methods:boxmodel}

We adopt the 5-box salinity model introduced by \citet{Wood20195BoxModel} and further reduced to 3 boxes by \citet{Alkhayuon2019AMOCRateTipping}. The original model partitions the Atlantic--Indo-Pacific basin into five boxes: North Atlantic ($N$), Tropical Atlantic ($T$), Southern Ocean ($S$), Indo-Pacific ($IP$), and a Deep box ($B$). In the implemented reduction, North and Tropical Atlantic salinity anomalies $S_N$ and $S_T$ are the two dynamical state variables. The slow variables of Southern Ocean and Deep-box salinities are prescribed as constants by the parameter vector, while Indo-Pacific salinity is derived at each time from salt conservation. The exact model equations and parameters are given in Appendix \ref{app:boxmodel}.

\citet{Wood20195BoxModel} calibrated two distinct parameter sets for the model: one for a pre-industrial climate ($p_{1\times}$) and one for a doubled-\CO{} climate ($p_{1\times}$) in the FAMOUS Earth System model. We construct a continuum of models through the full parameter vector as $p(t)=p_{1\times}+t(p_{2\times}-p_{1\times})$ and label it by $C(t)=280\,\text{ppm}+280\,\text{ppm}\,t$. Thus $t=0$ and $t=1$ correspond to 280 and 560 ppm, respectively, while $t>1$ denotes linear extrapolation beyond the doubled-\CO{} parameter vector. The nominal spacing is $\Delta t=0.05$, or $14\,\text{ppm}$, and we investigate the system from 280 through $896\,\text{ppm}$. We find that the on-state vanishes in a bifurcation when extrapolating the linear parameter ramp further, at around $900\,\text{ppm}$.

All numerical experiments and resilience calculations are performed using the Julia package Attractors.jl \citep{Datseris2023FrameworkGlobalStability} within the DynamicalSystems.jl ecosystem. This computational framework allows for the efficient integration of large initial condition ensembles and the simultaneous tracking of several resilience metrics along a parameter change \citep{Morr2026ComputingResilienceMeasures}. At each continuation step, we probe the two-dimensional state space with $10\,000$ initial conditions sampled uniformly from a $\pm 2\psu{}$ rectangle centred on the on-state attractor. Local resilience is computed from the Jacobian of the model equation's right-hand side. It is given as $-\mathrm{Re}(\lambda_\mathrm{max})$ where $\lambda_\mathrm{max}$ is the eigenvalue with the largest real part. In Figure \ref{fig:box-paper}, we map out the local basin geometry on a regular grid and show several example trajectories in terms of AMOC strength and phase-space evolution.


\begin{figure}[t]
  \centering
  \includegraphics[width=\textwidth,height=0.82\textheight,keepaspectratio]{amocbox_paper.png}
  \caption{\textbf{Three-box model response to salinity perturbations.}
  Top panels (a, b) show the evolution of the AMOC strength for four representative initial conditions, with dashed horizontal lines marking the on- and off-state equilibria. Bottom panels (c, d) show the same trajectories in the two-dimensional salinity state space (the two state variables), overlaid on the sampled basins of attraction; circles represent 0.05\psu{} neighbourhoods around the on- and off-state fixed points (blue and orange, respectively). The left and right columns correspond to 280 and $448\,\text{ppm}$ model specifications, respectively.}
  \label{fig:box-paper}
\end{figure}

\subsection{Boussinesq latitude-depth model}
\label{sec:methods:boussinesq}

The second model in our hierarchy is a latitude-depth Boussinesq model of the AMOC. The model evolves temperature, salinity, streamfunction, and vorticity on a $21\times41$ grid of discretised latitude and depth points. The one-dimensional sea surface is forced by cosine profiles of temperature and salinity flux, with lower temperatures and net freshwater input at high latitudes. This model is much less expensive than global circulation models but already contains spatially structured salinity fields and overturning, making it a useful intermediate step between the conceptual box model and the three-dimensional climate models.

We introduce a single external forcing parameter that mimics Arctic amplification: increasing it flattens the equator-to-pole surface temperature profile (Arctic amplification) and enhances the freshwater forcing profile (intensification of the hydrological cycle). For the exact mechanism of this parameter, its mapping to \CO{}, and the exact model equations, see Appendix \ref{app:boussinesq}. The perturbation experiments are performed across eight forcing values $\gamma=0.00,0.01,\ldots,0.07$, corresponding to \CO{} concentrations from $280\,\text{ppm}$ to approximately $455\,\text{ppm}$. For each configuration, the model is first spun up twice under weak hemispheric salinity asymmetries $\beta_s=0.1$ and $-0.1$. Under the symmetric salinity forcing $\beta_s=0$, for which an active and a reversed overturning state can coexist over part of the forcing range, we obtain the final AMOC on and off reference states by evolving the asymmetric-forcing equilibria until convergence. Based on the reference on-state, we initialise perturbed
198 salinity trajectories of the model. We then classify each perturbed run by the reference
state it settles to, which in turn defines the basin-volume and convergence-time metrics.

Initial-condition perturbations are applied to grid-aligned salinity boxes, and we utilise an affine conversion of the non-dimensional salinity units of the Boussinesq model to \psu{} (for details see Appendix \ref{app:boussinesq}). The North Atlantic box covers $40.5°$--$90°\text{N}$, the Tropical Atlantic box covers $31.5°\text{S}$--$40.5°\text{N}$, and the Southern compensation box covers the domain south of $31.5°\text{S}$. All three extend through the upper approximately $284\,\text{m}$ of the total $1000\,\text{m}$ water column. Cosine-tapered perturbation fields are normalised so that their mean over each hard box equals the requested amplitude. Perturbations are applied as salinity anomalies $(\Delta S_N,\Delta S_T)$ on a $20\times20$ grid ($400$ initial conditions), with both amplitudes spanning $\pm2\,\psu{}$ relative to the on-state equilibrium; a compensating perturbation in the Southern box conserves total domain salinity.


For the basin-mapping ensemble, only the salinity field of the on-state equilibrium is carried over into the perturbed initial conditions; the temperature, vorticity, and streamfunction fields are set to zero and left to re-adjust to the perturbed salinity under the model dynamics. This isolates the salinity perturbation as the driver of the response. To estimate local linear stability, we run a second, smaller ensemble of eight experiments per forcing level. Each run has a salinity perturbation in the two-dimensional (North and Tropical Atlantic) salinity space. The eight perturbations are evenly arranged on a circle of radius $\varepsilon = 0.5\,\psu{}$ around the on-state equilibrium. Unlike the basin-mapping ensemble, these perturbations are applied to the complete equilibrium state, retaining all temperature, vorticity, and streamfunction fields, so that the local recovery of the full system is reflected in the two-dimensional salinity projection. All trajectories are run under neutral symmetric salinity forcing for approximately $500$ model years. After discarding an initial ${\approx}50$-year transient, a $2\times 2$ discrete linear map $M$ is fitted by least squares to trajectory pairs $(x_t,\, x_{t+\Delta t})$ sampled every ${\approx}20$ years, where $x = (\delta S_N,\, \delta S_T)$ denotes the box-salinity deviation from equilibrium; trajectories that do not converge within the integration time are excluded. The continuous-time generator $A = \log(M)/\Delta t$ of the fitted map $M$ yields eigenvalues whose least stable real part gives the local resilience $-\mathrm{Re}(\lambda_{\max})$.

Throughout, we report salinities in \psu{} and the overturning strength in Sverdrups; the affine salinity calibration and linear overturning conversion, both anchored on the pre-industrial equilibrium, are given in Appendix \ref{app:boussinesq}. Time metrics are reported on a year-proxy scale equal to 100 times the model time unit. In Figure \ref{fig:boussinesq-paper}, we plot four exemplary trajectories in terms of the two box salinities and the AMOC strength at $280\,\text{ppm}$ and approximately $369\,\text{ppm}$.


\begin{figure}[t]
  \centering
  \includegraphics[width=\textwidth,height=0.82\textheight,keepaspectratio]{boussinesq_paper.png}
  \caption{\textbf{Boussinesq model response to salinity perturbations.}
  Top panels (a, b) show the evolution of the AMOC strength for four representative initial conditions, with dashed horizontal lines marking the on- and off-state equilibria. Bottom panels (c, d) show the same trajectories in North Atlantic and Tropical Atlantic box-mean salinity, overlaid on the sampled basins of attraction. White areas represent initial conditions that timed out or encountered negative salinity and are excluded from metrics. Circles represent 0.05\psu{} neighbourhoods around the on- and off-state fixed points (blue and orange, respectively). The left and right columns correspond to 280 and $369\,\text{ppm}$ model specifications, respectively.}
  \label{fig:boussinesq-paper}
\end{figure}


\subsection{CLIMBER-X}
\label{sec:methods:climberx}

CLIMBER-X \citep{Willeit2022EarthSystemModel} is a fast Earth system model of intermediate complexity (EMIC) that couples a 3-D ocean general circulation model, an energy-balance atmosphere, a thermodynamic sea-ice model, and land-surface components. The ocean component resolves the full three-dimensional salinity and temperature fields, including the main overturning circulation cells, at a resolution of approximately $5°\times 5°$ with 23 vertical levels. CLIMBER-X is computationally efficient enough to permit multi-millennial integrations, making it well-suited for exploring the state space through moderately sized perturbation ensembles.

We first equilibrate CLIMBER-X at 12 prescribed atmospheric \CO{} concentrations between 285 and $600\,\text{ppm}$: 285, 300, 330, 360, 400, 430, 460, 490, 510, 540, 570, and $600\,\text{ppm}$. Each equilibrium run is integrated for up to $20\,000$ model years from the default pre-industrial restart. After equilibration, the trailing $500$ years of the annual ocean time series are used to store reference values for AMOC strength at $26°\text{N}$ and for the salinity diagnostics in the geographical boxes illustrated in Fig.~\ref{fig:perturbations-overview}.

The North Atlantic box comprises Atlantic-basin grid cells from $35°$--$80°\text{N}$, the Tropical Atlantic box comprises Atlantic-basin cells from $35°\text{S}$--$35°\text{N}$, and the Southern compensation box spans all longitudes south of $35°\text{S}$. All boxes cover the upper $105\,\text{m}$, corresponding to the top four ocean layers. The North and Tropical box perturbations use a two-grid-cell latitudinal taper and a two-layer vertical cosine taper. The tapered fields are normalised so that the volume-weighted mean anomaly over the hard target box remains exactly equal to the prescribed salinity perturbation. A third compensating perturbation is applied in the Southern Ocean box so that the global ocean salt content in the restart is conserved. 

For our experiments, we apply instantaneous salinity perturbations to these boxes directly to the ocean restart as salinity anomalies in \psu{}, creating an ensemble of modified initial conditions. The North Atlantic and Tropical Atlantic perturbations each take the values $-2,-1.5,\ldots,2\psu{}$, giving a $9\times9$ grid of 81 experiments per \CO{} level. Each perturbed initial condition is integrated forward for $10\,000$ model years. This model is known to exhibit an AMOC off-state under high \CO{} concentrations when large constant freshwater forcing in the North Atlantic is applied. In the present setup of pure initial salinity perturbations, the off-state is not reached; instead, at higher \CO{} a second, distinct attractor with a somewhat stronger AMOC emerges. We classify each grid trajectory by its final state using the same criterion as the convergence-time reduction: a run is assigned to the on-state if its final $(s_{\mathrm{boxNA}},s_{\mathrm{boxTrop}})$ lies within an $\varepsilon=0.05\,\psu{}$ ball of the unperturbed reference state, and to the ``strong'' state otherwise. Figure \ref{fig:climberx-paper} shows four example trajectories starting from perturbed salinity initial conditions in the specified boxes, similar to the two figures above.

To estimate local linear stability, we run a second, smaller ensemble of $9$ experiments per \CO{} level. Eight runs apply isotropic salinity perturbations of magnitude $\varepsilon = 0.1\,\psu{}$ at angles $0°,45°,\ldots,315°$ in the two-dimensional (NA, Tropical) shallow-box salinity space, plus one unperturbed reference run; all perturbations use the same tapered, salt-conserving design as the grid ensemble. Each run is integrated for $500$ model years. An equilibrium reference is obtained from the mean of the unperturbed run over years 10--150. A $2\times2$ discrete linear map is fitted by least squares to annual trajectory pairs $(x_t,x_{t+1})$ in the same fitting window. The continuous-time generator $A=\log(M)/(1\,\text{yr})$ yields eigenvalues whose least stable real part gives the local resilience $-\mathrm{Re}(\lambda_{\max})$.


\begin{figure}[t]
  \centering
  \includegraphics[width=\textwidth,height=0.82\textheight,keepaspectratio]{climberx_paper.png}
  \caption{\textbf{CLIMBER-X model response to salinity perturbations.}
  Top panels (a, b) show the AMOC strength at $26°\text{N}$ along four example trajectories from perturbed salinity initial conditions. Bottom panels (c, d) show the same trajectories in the two-dimensional Atlantic-box salinity plane $(s_{\mathrm{boxNA}},s_{\mathrm{boxTrop}})$, with axes spanning exactly the $\pm2\,\psu{}$ perturbation window. The shaded background maps the basin of attraction of every initial condition on the full $9\times9$ grid: cells whose trajectory returns to the reference on-state (final salinity within $\varepsilon=0.05\,\psu{}$) are blue, and cells that settle on the distinct stronger-AMOC (``strong'') state are purple. Circles mark the $\varepsilon=0.05\,\psu{}$ neighbourhoods of the on-state (blue, at the unperturbed equilibrium) and of the strong state (purple, at the mean endpoint of the non-converged trajectories). The two columns show prescribed atmospheric \CO{} concentrations of $285\,\text{ppm}$ and $330\,\text{ppm}$: at $285\,\text{ppm}$ all 81 initial conditions return to the on-state, whereas at $330\,\text{ppm}$ 45 of the 81 converge to the strong state, which lies only ${\approx}0.05$--$0.08\,\psu{}$ from the on-state but exhibits an AMOC about $0.4\,\Sv{}$ stronger.}
  \label{fig:climberx-paper}
\end{figure}

\subsection{PlaSim-LSG}
\label{sec:methods:plasim}

The last component of the hierarchy is the coupled climate model of intermediate complexity PlaSim-LSG, comprising a spectral primitive-equation atmosphere (T21, $\approx 5.6°\times 5.6°$, 10 vertical levels) and the Large Scale Geostrophic (LSG) ocean \citep{Maier-Reimer1993MeanCirculationHamburg}, a full 3-D ocean circulation model at $3.5°\times 3.5°$ resolution with 22 vertical levels, coupled to a thermodynamic sea-ice scheme and a hydrological cycle.  With approximately $10^5$ degrees of freedom and a throughput of roughly 700 model years per CPU day, PlaSim-LSG is considerably more complex than CLIMBER-X. In particular, it resolves atmospheric variability more faithfully and thereby yields seemingly chaotic ocean dynamics.

At $285\,\text{ppm}$ (pre-industrial) and $360\,\text{ppm}$ (current) \CO{}, PlaSim-LSG exhibits a bistable AMOC with an on-state of $\approx 15\Sv{}$ and an off-state (deep convection in Subpolar Gyre ceases) of $\approx 1-2\Sv{}$, both confirmed by long unforced simulations. \citet{Borner2025AMOCGlobalStability} applied the edge-tracking algorithm \citep{Lucarini2017EdgeTracking} to locate the edge state (Melancholia state) of the AMOC in PlaSim-LSG at $285\,\text{ppm}$ and $360\,\text{ppm}$.  The algorithm iteratively bisects pairs of trajectories converging to opposite attractors and tracks the resulting pseudotrajectory along the basin boundary until it converges to the edge state. The edge state is a chaotic saddle characterized by persistent centennial AMOC oscillations ($\approx 120$\,yr period, amplitude $3$–$10\Sv{}$), driven by atmosphere–ice–ocean interactions in the North Atlantic. A complete description of the model configuration, its AMOC bistability, and the edge-tracking procedure is given in \citet{Borner2025AMOCGlobalStability}, to which we refer the reader for full methodological details.

From the final iterations of edge tracking at each \CO{} level, we extract pairs of edge tracks: one converging to the on-state and one converging to the off-state, both initialised from near-boundary conditions that closely shadow the edge state before diverging. These trajectories provide a natural set of initial conditions from which resilience-relevant transient behaviour can be quantified: they start close to the basin boundary and thus represent challenging recovery paths for the on-state. In total, we analyse 37 paired edge-track files at $285\,\text{ppm}$ and 38 at $360\,\text{ppm}$, with two trajectories per file. Tracks are classified from their final North Atlantic salinity, with the higher-salinity endpoint identified as AMOC-on, and convergence statistics include only trajectories classified as converged. Dedicated long equilibrium simulations of the AMOC-on, AMOC-off, and edge states at each \CO{} level provide the attractor clouds against which these transients are evaluated.

We reduce the high-dimensional PlaSim-LSG salinity output to the salinity of two geographical boxes. The North Atlantic readout is the mean salinity from $35°$--$80°\text{N}$ and 0--$1000\,\text{m}$ depth, and the Tropical Atlantic readout covers $35°\text{S}$--$35°\text{N}$ and 0--$500\,\text{m}$ depth \todo{check depth}. For the edge tracks, these averages are applied to Atlantic zonal-mean salinity; for the equilibrium runs, an Atlantic basin mask is applied before zonal averaging. Dedicated on-, off-, and edge-state equilirbium trajectories determine the means associated with these states in this two-dimensional salinity plane. The covariance matrices of the scattered trajectory points are computed to obtain an approximate Gaussian distribution of data for each state. Convergence time is measured from the last visit to the edge-state $4\sigma$ Gaussian ellipse to the first entry into the target-attractor ellipse, and the critical-shock proxy is the gap between the edge-state and on-state ellipses. Local resilience is the covariance proxy $[2000\sqrt{\lambda_{\max}(C)}]^{-1}$ for the on-state equilibrium covariance $C$ \todo{width of the ellipse}.

\begin{figure}[t]
  \centering
  \includegraphics[width=\textwidth,height=0.82\textheight,keepaspectratio]{plasim_paper.png}
  \caption{\textbf{PlaSim-LSG edge-state trajectories and covariance ellipses of attractors.}
  Top panels (a, b) show AMOC strength readout from the PlaSim model along filtered edge-track trajectories. Bottom panels (c, d) show the same transients in the leading reduced salinity coordinates, together with equilibrium-run clouds for the on-, off-, and edge-states (trajectories not shown). Ellipses denote the fitted $3\sigma$ Gaussian variance used to define convergence times and edge-to-attractor distances. The two columns show the experiment results at prescribed atmospheric \CO{} concentration at $285\,\text{ppm}$ and $360\,\text{ppm}$, respectively.}
  \label{fig:plasim-paper}
\end{figure}

\section{Results}
\label{sec:results}

\begin{figure}[p]
  \centering
  \includegraphics[width=\textwidth,height=0.82\textheight,keepaspectratio]{synthesis_figure.png}
  \caption{\textbf{Cross-model synthesis of AMOC resilience metrics.}
  Resilience quantifiers are plotted against the \CO{} coordinate for each model, where available. The model-specific definitions and implementation of each resilience quantifier are given in Table \ref{tab:overview}. These results provide the basis for assessing which resilience signals are coherent across the model hierarchy.}
  \label{fig:synthesis}
\end{figure}
\FloatBarrier

\section{Discussion}
\label{sec:discussion}

\todo{Coherence of resilience metrics across the hierarchy, Global versus local metrics, Implications for early-warning detection, Limitations and caveats}



\appendix

\section{Specifics on the 3-Box salinity model}
\label{app:boxmodel}

The model tracks North Atlantic and Tropical Atlantic salinity anomalies $S_N$ and $S_T$ in the internal convention $S_i=100(S_i^{\mathrm{dim}}-S_0)$, where $S_0=0.035$. In the figures, these values are converted into practical salinity as $S_i^{\mathrm{plot}}=35+10S_i\psu{}$. Southern Ocean salinity $S_S$ and Deep-box salinity $S_B$ are prescribed and converted to the same anomaly convention. Indo-Pacific salinity is instead diagnosed from total salt content $C$ at every model evaluation:
\begin{equation}
S_{IP}=\frac{100}{V_{IP}}\left[C-\frac{V_NS_N+V_TS_T+V_SS_S+V_BS_B}{100}-S_0(V_N+V_T+V_S+V_{IP}+V_B)\right].
\label{eq:sip}
\end{equation}

The hosing-adjusted freshwater fluxes are $\widetilde F_N=F_N+0.1311\times10^6H$ and $\widetilde F_T=F_T+0.6961\times10^6H$. The continuation used here has $H=0$ in both endpoint parameter vectors. The implemented governing equations are
\begin{align}
\frac{\dint S_N}{\dint t}
&=\begin{cases}
\dfrac{Y}{V_N}\left[\dfrac{(q+K_N)(S_T-S_N)}{100}-\widetilde F_NS_0\right], & q\geq0,\\[6pt]
\dfrac{Y}{V_N}\left[-\dfrac{q(S_B-S_N)}{100}+\dfrac{K_N(S_T-S_N)}{100}-\widetilde F_NS_0\right], & q<0,
\end{cases}
\label{eq: dS_N}\\[4pt]
\frac{\dint S_T}{\dint t}
&=\begin{cases}
\begin{aligned}
\dfrac{Y}{V_T}\Bigl[&\dfrac{q(\gamma S_S+(1-\gamma)S_{IP}-S_T)}{100}+\dfrac{K_S(S_S-S_T)}{100}\\[-2pt]
&+\dfrac{K_N(S_N-S_T)}{100}-\widetilde F_TS_0\Bigr],
\end{aligned} & q\geq0,\\[10pt]
\begin{aligned}
\dfrac{Y}{V_T}\Bigl[&-\dfrac{q(S_N-S_T)}{100}+\dfrac{K_S(S_S-S_T)}{100}\\[-2pt]
&+\dfrac{K_N(S_N-S_T)}{100}-\widetilde F_TS_0\Bigr],
\end{aligned} & q<0,
\end{cases}
\label{eq: dS_T}\\[4pt]
q&=\frac{\lambda\left[\alpha(T_S-T_0)+\beta(S_N-S_S)/100\right]}{1+\lambda\alpha\mu}.
\label{eq: q}
\end{align}
Here $V_i$ are box volumes, $K_N$ and $K_S$ are mixing coefficients, $F_N$ and $F_T$ are freshwater fluxes, and $Y$ is the model's time-scaling parameter. The routing parameter $\gamma$ sets the fraction of positive overturning return flow passing through the Southern Ocean rather than the Indo-Pacific. The overturning flow $q$ is reported in $\text{m}^3\,\text{s}^{-1}$ and converted to Sverdrups by division by $10^6$. Its dependence on flow direction makes the right-hand side discontinuous at $q=0$.

The parameter vectors are taken from \citet{Wood20195BoxModel} and \citet{Alkhayuon2019AMOCRateTipping}, who calibrated pre-industrial ($1\times$\CO{}) and doubled-\CO{} ($2\times$\CO{}) configurations against the FAMOUS Earth System model. We linearly interpolate and extrapolate the full vectors, including changing box volumes, freshwater fluxes, temperatures, mixing coefficients, hydraulic and routing parameters, and total salt content, as described in Section~\ref{sec:methods:boxmodel}.

\section{Specifics on the Boussinesq model}
\label{app:boussinesq}

\subsection*{Physical setup}

The model represents a two-dimensional meridional cross-section of the ocean at non-dimensional latitude $x \in [0, 5]$ and depth $z \in [0, 1]$, where $x = 0$ corresponds to the South Pole, $x = 2.5$ to the equator, and $x = 5$ to the North Pole; $z = 0$ is the ocean floor and $z = 1$ is the surface \citep{Esclapez2026ConstructingEfficientScore}. The zonally averaged flow is spatially two-dimensional and incompressible under the Boussinesq approximation, in which density variations are retained only in the buoyancy term. The velocity field is derived from a streamfunction $\psi(x,z,t)$ via
\begin{equation}
  u = \frac{\partial \psi}{\partial z}, \qquad w = -\frac{\partial \psi}{\partial x},
  \label{eq:vel_bous}
\end{equation}
where $u$ is the horizontal (meridional) velocity and $w$ is the vertical velocity. The vorticity is $\omega = -\nabla^2\psi$. The non-dimensional base parameters and numerical resolution are given in Table~\ref{tab:boussinesq-params}; the salinity asymmetry is changed between spin-up and perturbed integrations as described below.

\begin{table}[h]
  \centering
  \caption{Default non-dimensional parameters for the Boussinesq model.}
  \label{tab:boussinesq-params}
  \begin{tabular}{lll}
    \toprule
    Symbol & Value & Physical meaning \\
    \midrule
    $\mathrm{Pr}$   & $1$              & Prandtl number ($\nu/\kappa_T$) \\
    $\mathrm{Le}$   & $1$              & Lewis number ($\kappa_T/\kappa_S$) \\
    $\mathrm{Ra}$   & $4\times10^4$    & Rayleigh number \\
    $a$             & $5$              & Domain length (pole-to-pole) \\
    $\delta$        & $0.05$           & Surface boundary-layer thickness \\
    $\tau_T$        & $0.1$            & Thermal restoring time scale \\
    $\tau_S$        & $1.0$            & Salinity flux time scale \\
    $\beta_s$       & $0.1$            & Spin-up salinity hemispheric asymmetry \\
    $M,\,N$         & $20,\,40$        & Grid intervals ($21\times41$ points) \\
    $\Delta t$      & $10^{-3}$        & Time step \\
    \bottomrule
  \end{tabular}
\end{table}

\subsection*{Governing equations}

Temperature $T$, salinity $S$, and vorticity $\omega$ evolve according to coupled advection–diffusion equations with surface forcing:
\begin{align}
  \frac{\partial T}{\partial t} + u\frac{\partial T}{\partial x} + w\frac{\partial T}{\partial z}
    &= \nabla^2 T + \frac{h(z)}{\tau_T}\bigl(T_s^\gamma(x) - T\bigr), \label{eq:bous_T}\\[4pt]
  \frac{\partial S}{\partial t} + u\frac{\partial S}{\partial x} + w\frac{\partial S}{\partial z}
    &= \frac{1}{\mathrm{Le}}\nabla^2 S + \frac{h(z)}{\tau_S}\,S_s^\gamma(x), \label{eq:bous_S}\\[4pt]
  \frac{\partial \omega}{\partial t} + u\frac{\partial \omega}{\partial x} + w\frac{\partial \omega}{\partial z}
    &= \mathrm{Pr}\,\nabla^2\omega + \mathrm{Pr}\cdot\mathrm{Ra}\,\frac{\partial}{\partial x}(T - S), \label{eq:bous_omega}\\[4pt]
  \nabla^2\psi &= -\omega. \label{eq:bous_psi}
\end{align}

The surface boundary-layer profile $h(z) = \exp(-(1-z)/\delta)$ confines the surface forcing to the uppermost ocean. The temperature equation relaxes towards the prescribed surface profile with time scale $\tau_T$; a smaller $\tau_T$ enforces temperature more strongly towards $T_s$. For salinity, a fixed flux formulation is used instead of a relaxation, which is physically appropriate since ocean salinity is set by freshwater exchange with the atmosphere rather than by direct relaxation to a reference profile. A global salinity correction is applied at each time step to conserve domain-mean salinity exactly.

The surface temperature profile is
\begin{equation}
  T_s(x) = \tfrac{1}{2}\!\left[\cos\!\left(\tfrac{2\pi x}{a} - \pi\right) + 1\right],
  \label{eq:Ts}
\end{equation}
giving $T_s = 0$ at both poles and $T_s = 1$ at the equator. This equator-to-pole contrast drives the horizontal density gradient that sustains the overturning. The surface salinity flux is
\begin{equation}
  S_s(x) = 3.5\cos\!\left(\tfrac{2\pi x}{a} - \pi\right) - \beta_s\sin\!\left(\tfrac{\pi x}{a} - \tfrac{\pi}{2}\right),
  \label{eq:Ss}
\end{equation}
where the dominant cosine term represents the atmospheric evaporation–precipitation cycle (net evaporation in the tropics, net precipitation at high latitudes) and the $\beta_s$ term introduces a weak hemispheric asymmetry that adds a slight freshwater excess in the northern high latitudes relative to the south.

The vorticity equation~\eqref{eq:bous_omega} is obtained by taking the curl of the momentum equation. The baroclinic production term $\mathrm{Pr}\cdot\mathrm{Ra}\,\partial_x(T-S)$ encodes the linearised equation of state under the Boussinesq approximation: density anomalies $\rho \propto -(T - S)$ create horizontal density contrasts that drive overturning vorticity. The product $\mathrm{Pr}\cdot\mathrm{Ra} = 4\times10^4$ governs the balance between advective and diffusive vorticity transport and is the primary control on whether a sustained circulation develops.

\subsection*{Boundary conditions and numerical discretisation}

No-flux ($\partial_n T = \partial_n S = 0$) and no-slip conditions are imposed on all lateral boundaries. The streamfunction satisfies $\psi = 0$ on all walls, consistent with a closed domain with no net overturning flux through any boundary. The depth coordinate is discretised on a non-uniform grid that concentrates points near the surface and the ocean floor via the mapping
\begin{equation}
  z_j = \tfrac{1}{2} + \frac{\tanh\!\bigl(3(j/N - \tfrac{1}{2})\bigr)}{2\tanh(3/2)}, \qquad j = 0,\ldots,N,
  \label{eq:zgrid}
\end{equation}
and the latitude coordinate uses a uniform grid. Spatial derivatives are approximated by second-order finite differences. Equations~\eqref{eq:bous_T}--\eqref{eq:bous_omega} are integrated using an implicit-explicit (IMEX) scheme with time step $\Delta t = 10^{-3}$: diffusion and surface-restoring terms are treated implicitly to avoid the diffusive CFL restriction, while advection terms are treated explicitly. The implicit step at each time leads to a separable Sylvester equation that is solved with the Bartels--Stewart algorithm.

\subsection*{Arctic amplification and \CO{} mapping}

To sweep the model across different climate states, the surface temperature forcing is modified to mimic polar amplification by blending the standard profile towards its spatial mean:
\begin{equation}
  T_s^\gamma(x) = (1-\gamma)\,T_s(x) + \gamma\overline{T_s},
  \label{eq:Tgamma}
\end{equation}
where $\overline{T_s}$ is the mean of the unmodified temperature profile on the model grid and $\gamma \in [0,1]$ measures the degree of Arctic amplification. At $\gamma = 0$ the forcing is unchanged (pre-industrial); increasing $\gamma$ reduces the equator-to-pole temperature gradient, weakening the baroclinic torque in \eqref{eq:bous_omega} and thereby reducing the thermal driving of the overturning.

The same parameter also intensifies the hydrological forcing: under a warming climate the evaporation--precipitation contrast strengthens (``wet-gets-wetter''), which we represent by scaling the symmetric component of the surface salinity flux \eqref{eq:Ss} while leaving the hemispheric asymmetry unchanged,
\begin{equation}
  S_s^\gamma(x) = (1 + \alpha_s\gamma)\,3.5\cos\!\left(\tfrac{2\pi x}{a} - \pi\right) - \beta_s\sin\!\left(\tfrac{\pi x}{a} - \tfrac{\pi}{2}\right),
  \label{eq:Sgamma}
\end{equation}
with hydrological-intensification coefficient $\alpha_s = 1$, so that the evaporation--precipitation amplitude would double at $\gamma = 1$ and increases by up to $7\%$ across the sweep ($\gamma \le 0.07$). Increasing $\gamma$ thus simultaneously weakens the thermal driving and strengthens the salinity forcing of the overturning.

A physically grounded mapping between $\gamma$ and atmospheric \CO{} concentration $C$ is derived from the logarithmic forcing formula and an estimate of polar ocean-surface amplification:
\begin{equation}
  \gamma(C) = \frac{(\alpha_\text{sst}-1)\,\text{ECS}}{\Delta T_0}\log_2\!\left(\frac{C}{C_0}\right)
  = 0.1\log_2\!\left(\frac{C}{280\,\text{ppm}}\right).
  \label{eq:co2_gamma}
\end{equation}
We adopt the IPCC AR6 best-estimate equilibrium climate sensitivity $\text{ECS} = 3.0\,^\circ\text{C}$ per \CO{} doubling \citep{Forster2021IPCCAR6Ch7}, a polar ocean-surface amplification factor $\alpha_\text{sst} = 2.0$ consistent with observed Arctic SST warming being approximately two- to threefold the global-mean increase, and a pre-industrial Atlantic equator-to-pole SST difference $\Delta T_0 = 30\,^\circ\text{C}$ (approximate Levitus climatology value). These choices give $\gamma = 0$ at pre-industrial \CO{} ($C_0 = 280\,\text{ppm}$) and $\gamma = 0.1$ at doubled \CO{} ($560\,\text{ppm}$). The parameter sweep covers eight uniformly spaced values of $\gamma \in [0,0.07]$ (steps of $0.01$), corresponding to \CO{} concentrations from $280\,\text{ppm}$ to approximately $455\,\text{ppm}$.

\subsection*{Perturbation geometry and physical scaling}

Each salinity box is a finite set of grid cells whose lateral and lower boundaries are snapped to grid-cell faces. The North Atlantic box covers $x\in[3.625,5]$ (approximately $40.5^\circ$--$90^\circ\text{N}$), the Tropical Atlantic box covers $x\in[1.625,3.625]$ (approximately $31.5^\circ\text{S}$--$40.5^\circ\text{N}$), and the Southern compensation box covers the domain south of $x=1.625$. All boxes include cells with $z\geq1-0.28449$, corresponding to the top 15 levels and a lower face near $284\,\text{m}$ on the $1000\,\text{m}$ plotting scale. A cosine taper is applied at interior box edges while domain-wall edges remain sharp. Each tapered field is normalised so that its mean over the corresponding hard membership mask equals one; adding amplitude $\Delta S$ therefore changes the box-mean salinity by exactly $\Delta S$. North Atlantic and Tropical Atlantic perturbations are compensated in the Southern box to preserve the model's domain cell-mean salinity.

A spin-up with weak hemispheric asymmetry $\beta_s=0.1$ seeds the symmetric-forcing ($\beta_s=0$) AMOC-on equilibrium. A second $\beta_s=-0.1$ spin-up similarly seeds the reversed-AMOC equilibrium. At the start of each ensemble, both spin-up states are released under $\beta_s=0$ dynamics and integrated until they settle. Basin-sweep initial conditions retain only the perturbed salinity field of the $\beta_s=0.1$ spin-up; temperature, vorticity, and streamfunction are set to zero. A trajectory is stopped when its box salinities enter the $\varepsilon=0.05\,\psu{}$ neighbourhood of either startup equilibrium. Valid final streamfunction values are clustered independently, the strongest cluster is designated AMOC-on, and basin volume is the fraction assigned to that cluster. Mean convergence time is then evaluated for those AMOC-on trajectories relative to the median final box salinities of the cluster.

Model salinities are mapped to practical salinity units by the affine calibration $S\,[\psu{}]=10.72735\,s+25.45555$, fixed by requiring the North Atlantic and Tropical boxes of the pre-industrial ($\gamma=0$, $\beta_s=0.1$) equilibrium to match $35.0$ and $36.5\psu{}$; the $\pm2\psu{}$ perturbation grid therefore corresponds to approximately $\pm0.18644$ model units. Likewise, the overturning streamfunction is converted to Sverdrups by anchoring the pre-industrial equilibrium value $\psi_0\approx1.517$ model units to $18\Sv{}$, giving a factor of approximately $11.864\Sv{}/\psi_\text{model}$. Convergence and characteristic return times are multiplied by 100 to obtain the year-proxy scale used in the synthesis, and local resilience is divided by 100 to obtain inverse year-proxy units.

\section{Specifics on the CLIMBER-X model}
\label{app:regions}

\section{Specifics on the PlaSim-LSG model}
\label{app:eof}



\bibliographystyle{apalike}
\bibliography{BibAll}


\end{document}
